\documentclass[11pt]{article}

\usepackage[margin=1in]{geometry}
\usepackage{amsmath,amssymb}
\usepackage{booktabs}
\usepackage{multirow}
\usepackage{graphicx}
\usepackage{xcolor}
\usepackage{enumitem}
\usepackage{caption}
\usepackage[hidelinks]{hyperref}
\usepackage{microtype}
\usepackage{authblk}

\graphicspath{{../figures/}}

\title{\textbf{OncoEvidence: when does a knowledge graph earn its place\\
in oncology drug repurposing?}\\[2pt]
\large A mechanism-first, honestly-evaluated evidence-triage pipeline}

\author{OncoEvidence}
\affil{\small Research preprint --- hypothesis-generating, not medical advice}
\date{}

\begin{document}
\maketitle

\begin{abstract}
\noindent
Knowledge-graph (KG) neural networks are widely promoted for drug repurposing, usually on
the claim that the graph improves link prediction. We test that claim honestly on PrimeKG
oncology indications and find it does \emph{not} hold: a properly tuned tabular model
(XGBoost) on the identical node features matches or beats a heterogeneous graph neural
network (GNN) in every evaluation regime, including inductive cold-disease and cold-drug
splits. Rather than discard the graph, we move its job to where a tabular model cannot
follow --- producing a \emph{traceable multi-hop mechanism-of-action (MOA) path}
($\text{drug}\rightarrow\text{target}\rightarrow(\text{PPI}/\text{pathway})\rightarrow\text{cancer
gene}\rightarrow\text{cancer}$) --- and we add a literature-grounded LLM layer that checks
whether each proposed path is supported by the published record. We evaluate this pipeline
instead of demonstrating it. Mechanism structure separates true indications from random
pairs at AUROC $0.879$ and agrees with curated DrugMechDB mechanisms at $0.802$; under
deliberately hard negatives the signal is strong against random and oncology-drug controls
($0.887$/$0.870$) but only modest ($0.609$) against shared-target negatives, which we report
rather than hide. A prospective time-split shows the graph \emph{enriches a small,
top-priority tail}: trained only on pre-2005 structure it surfaces $11\times$ as many
future indications in the top 50 as a structure-blind control (recovering post-cutoff
approvals such as romidepsin and pralatrexate for cutaneous T-cell lymphoma and nilotinib
for CML), although it is not better on median rank. A counterfactual test confirms the
mechanism head is faithful: deleting the single curated mechanism edge collapses the
prediction far more than deleting a random edge ($p\approx3.6\times10^{-34}$). Finally, a
popularity-deconfounded ClinicalTrials.gov check is reported as a fully characterized
negative. The contribution is not a new model but an honest account of \emph{when graph
structure is actually necessary} and a mechanism-first pipeline that acts on that finding.
\end{abstract}

\section{Introduction}
Drug repurposing --- finding new oncology uses for existing drugs --- is attractive because
the safety profile is known and development is cheaper than \emph{de novo} discovery.
Biomedical knowledge graphs such as PrimeKG~\cite{primekg} encode drugs, diseases, proteins,
and pathways as a heterogeneous graph, and a large literature applies graph neural networks
(GNNs) to predict drug--disease links over them~\cite{txgnn,decagon}. The usual headline is
that the graph improves prediction.

We argue that this headline is often an artifact of a missing baseline. When the same node
features that feed the GNN are handed to a \emph{tuned} gradient-boosted tree, the tabular
model is an extremely strong ranker, and the relational structure adds little to the link
task. Reporting this negative result is the honest starting point of our work, and it
reframes the question from ``does the graph win?'' to ``\emph{what is the graph uniquely
good for?}''

Our answer: the graph's irreplaceable output is not a score but a \emph{mechanism}. A
tabular model can rank a (drug, cancer) pair; it cannot tell you \emph{why}. The KG can
produce a short, readable, biologically-typed chain from the drug, through its protein
target and interaction/pathway neighborhood, to a cancer-associated gene and the cancer.
We make that chain the deliverable, filter out phenotype/symptom coincidences, and add an
LLM evidence layer that grounds each chain in retrieved literature. We then \emph{evaluate}
the whole pipeline against curated mechanisms, hard negatives, a prospective time-split, and
a counterfactual faithfulness test.

\paragraph{Contributions.}
\begin{enumerate}[leftmargin=1.5em,itemsep=2pt,topsep=2pt]
\item \textbf{An honest ``is the graph necessary?'' audit.} With a tuned XGBoost on
identical features (plus topology-removal and memorization-control ablations), the GNN does
not beat the tabular baseline on link ranking in any regime (Sec.~\ref{sec:link}).
\item \textbf{Mechanism-path extraction as the graph's real job}, prioritizing MOA relations
over phenotype/symptom overlap, with a quantitative true-vs-random separation and an
\emph{honest hard-negative stress test} (Sec.~\ref{sec:mech}).
\item \textbf{A literature-grounded LLM evidence layer} that reads abstracts (not just
titles) and is more precise than a lexical baseline against curated mechanisms
(Sec.~\ref{sec:llm}).
\item \textbf{Prospective and faithfulness audits}: top-tail enrichment of future
indications under a time-split (Sec.~\ref{sec:prospective}) and a counterfactual
edge-deletion test of mechanism faithfulness (Sec.~\ref{sec:faith}).
\item \textbf{A popularity-deconfounded orthogonal check} against ClinicalTrials.gov,
reported transparently as a negative (Sec.~\ref{sec:ct}).
\end{enumerate}
We do not claim to invent KG- or LLM-based repurposing; close prior work
exists~\cite{txgnn,kgmlxdtd}. The contribution is the honest evaluation and the
mechanism-first design that follows from it.

\section{Data and models}
We build a heterogeneous graph from PrimeKG~\cite{primekg} with node types for drugs,
diseases, genes/proteins, pathways, and biological processes, and the corresponding typed
relations (indication, drug--protein, disease--protein, protein--protein interaction,
pathway membership). Node features are sentence-embedding encodings of node names/descriptions
(with a hashing fallback for offline use). Oncology diseases are flagged via an ontology
filter. The target task is the drug--disease \texttt{indication} relation.

We compare four models under a shared evaluation harness: (i) a 2-layer heterogeneous GNN
with an InfoNCE-style link decoder; (ii) a tuned XGBoost on concatenated endpoint features;
(iii) a feature MLP; and (iv) a DistMult KG embedding as a memorization control. We evaluate
three regimes: \emph{transductive}, \emph{inductive cold-disease} (held-out oncology diseases
never seen in training), and \emph{inductive cold-drug}. Negatives are sampled with degree
awareness and all splits are leakage-checked.

\section{The link task does not need the graph}\label{sec:link}
Table~\ref{tab:bench} reports the four-model comparison over five seeds. In the transductive
regime the tuned XGBoost leads (AUROC $0.988$ vs the GNN's $0.977$). The gap matters most in
the inductive cold-disease oncology regime --- the setting closest to real repurposing,
where the disease is unseen --- where XGBoost reaches $0.963$ against the GNN's $0.882$. The
GNN never wins. A topology ablation makes the reason concrete: shuffling the edges barely
changes GNN performance, so it is leaning on node features rather than relational structure
for this task.

\begin{table}[t]
\centering\small
\caption{Link-prediction benchmark (mean$\pm$sd over 5 seeds). The tuned tabular model
matches or beats the GNN in every regime; the graph does not earn its place on link ranking
alone.}
\label{tab:bench}
\begin{tabular}{llccc}
\toprule
\textbf{Regime} & \textbf{Model} & \textbf{AUROC} & \textbf{AUPRC} & \textbf{F1}\\
\midrule
\multirow{2}{*}{Transductive} & XGBoost & \textbf{0.988$\pm$0.002} & \textbf{0.988$\pm$0.002} & \textbf{0.952$\pm$0.003}\\
 & GNN & 0.977$\pm$0.003 & 0.972$\pm$0.005 & 0.937$\pm$0.005\\
\midrule
\multirow{2}{*}{Cold-disease (onc.)} & XGBoost & \textbf{0.963$\pm$0.007} & \textbf{0.954$\pm$0.013} & \textbf{0.919$\pm$0.007}\\
 & GNN & 0.882$\pm$0.050 & 0.828$\pm$0.052 & 0.843$\pm$0.040\\
\midrule
\multirow{2}{*}{Cold-drug} & XGBoost & \textbf{0.958$\pm$0.007} & \textbf{0.957$\pm$0.007} & 0.906$\pm$0.010\\
 & GNN & 0.956$\pm$0.006 & 0.943$\pm$0.008 & \textbf{0.933$\pm$0.008}\\
\bottomrule
\end{tabular}
\end{table}

\section{The graph's real job: mechanism paths}\label{sec:mech}
We connect a drug to a disease through \emph{mechanism} relations only, in three templates of
decreasing specificity:
\[
\begin{aligned}
\text{direct target:}\quad & \text{drug}\xrightarrow{\text{targets}} P \xleftarrow{\text{assoc}} \text{disease}\\
\text{PPI bridge:}\quad & \text{drug}\xrightarrow{\text{targets}} P_1 \xrightarrow{\text{PPI}} P_2 \xleftarrow{\text{assoc}} \text{disease}\\
\text{shared pathway:}\quad & \text{drug}\xrightarrow{\text{targets}} P_1 \xrightarrow{\text{in}} \text{pw} \xleftarrow{\text{in}} P_2 \xleftarrow{\text{assoc}} \text{disease}
\end{aligned}
\]
Phenotype/symptom relations (a drug whose side effect happens to be a disease symptom) are
\emph{excluded} by construction, and promiscuous hub intermediates (plasma carriers such as
albumin, generic high-degree proteins) are blocked or down-weighted by an IDF-style
promiscuity penalty so the path stays specific. Each pair gets a single mechanism score equal
to its strongest path.

\paragraph{Does mechanism structure separate true from random pairs?} Over 400 true
oncology indications and 400 random drug--cancer pairs, the mechanism score separates the
groups at \textbf{AUROC 0.879}, with a direct-target rate of $34\%$ for true pairs versus
$0.25\%$ for random, and the extracted chains recover textbook MOAs (FLT3, RARA, TOP2A,
TYMS).

\paragraph{Honest hard negatives.} A single number against random pairs is too easy, so we
stress-test against progressively harder negatives (Table~\ref{tab:hardneg}). The signal
stays strong against random and oncology-drug negatives ($0.887$/$0.870$) and degrades
gracefully against degree-matched negatives ($0.742$), but against \emph{shared-target}
negatives --- a drug whose target is itself linked to the cancer --- it falls to $0.609$,
near chance. This is the honest boundary of the method: mechanism structure is decisive
against generic negatives and only weakly discriminative when the negative already shares the
same biology. We foreground this rather than reporting only the favorable number.

\begin{table}[t]
\centering\small
\caption{Hard-negative stress test (separation AUROC, true vs.\ each negative class). Strong
vs.\ random/oncology-drug; modest under the hardest shared-target controls.}
\label{tab:hardneg}
\begin{tabular}{lcccc}
\toprule
\textbf{Negative class} & \textbf{random} & \textbf{oncology-drug} & \textbf{degree-matched} & \textbf{shared-target}\\
\midrule
Separation AUROC & 0.887 & 0.870 & 0.742 & \textbf{0.609}\\
\bottomrule
\end{tabular}
\end{table}

\paragraph{Curated agreement.} Mapped to HGNC symbols, the extracted mechanism genes agree
with expert-curated DrugMechDB mechanisms at \textbf{0.802} on covered pairs --- the paths
are not just internally consistent, they match an external ground truth.

\section{Literature-grounded LLM verification}\label{sec:llm}
For each candidate the pipeline retrieves Europe PMC records (title \emph{and abstract}) and
assembles the mechanism path plus the retrieved abstracts into a retrieval-augmented prompt.
An LLM writes a structured evidence dossier (mechanistic rationale / supporting /
contradicting / verdict) grounded in the supplied abstracts, and a separate LLM-as-judge
pass scores plausibility and evidence strength. Against the curated mechanisms, the LLM's
\emph{supported} calls are precise (precision $0.857$) versus $0.591$ for a lexical
baseline, and a stricter MOA rubric with sentence-level grounding rejects coincidental hub
paths (e.g.\ albumin bridges). The LLM verifier grades true pairs \emph{supported/weak} and
random pairs \emph{no-path} ($11/50$ vs $1/50$ supported).

\section{Prospective audit: top-tail enrichment over time}\label{sec:prospective}
PrimeKG has no timestamps, so we assign each true indication the year of its earliest Europe
PMC co-mention and hold out indications established after a cutoff $T=2005$ ($350$ sampled
pairs, $341$ years resolved; $240$ past / $101$ future; 3 seeds). Trained only on pre-$T$
structure, the graph model ranks held-out future indications at \textbf{AUROC 0.930} (vs
sampled negatives), beating a structure-blind control by $+0.047$.

The aggregate ranking story is deliberately nuanced. When we instead rank every drug against
each future cancer using only pre-$T$ structure, the GNN's \emph{median} rank ($3{,}019$ of
$\sim$7{,}956) is \emph{worse} than the control's ($1{,}482$), and a paired Wilcoxon test
does not favor the GNN on overall rank ($p\approx1.0$). What the graph does is
\emph{enrich the top-priority tail}: it places \textbf{11 of 101} future indications in the
top 50 (a $17.3\times$ enrichment over a random screen) versus \textbf{1} for the control,
and $13.9\%$ in the top $\sim$1\% versus $1.0\%$. The hits are recognizable post-cutoff
approvals the model had no indication edge for (Table~\ref{tab:cases}). The honest claim is
narrow: the graph does not rank future indications better on average; it concentrates a real
signal exactly where a repurposing shortlist is read.

\begin{table}[t]
\centering\small
\caption{Prospective named cases ($T=2005$, single seed). A model trained on pre-$T$
structure ranks each drug near the top of $\sim$7{,}956 candidates for a cancer whose
indication is first co-mentioned only after $T$, far above the structure-blind control.}
\label{tab:cases}
\begin{tabular}{llccc}
\toprule
\textbf{Drug} & \textbf{Cancer (future indication)} & \textbf{1st lit.\ yr} & \textbf{GNN rank} & \textbf{MLP rank}\\
\midrule
Trabectedin & ovarian carcinoma (subtypes) & 2006--10 & \textbf{2--5} & 3038\\
Methoxsalen & cutaneous T-cell lymphoma & 2006 & \textbf{10} & 2154\\
Romidepsin & cutaneous T-cell lymphoma & 2008--10 & \textbf{63} & 1726\\
Pralatrexate & cutaneous T-cell lymphoma & 2009 & \textbf{321} & 2235\\
Nilotinib & CML, BCR-ABL1 positive & 2006 & \textbf{342} & 1098\\
Plerixafor & plasma cell myeloma & 2007 & \textbf{428} & 1797\\
\bottomrule
\end{tabular}
\end{table}

\section{Faithfulness: is the named mechanism causal to the prediction?}\label{sec:faith}
A mechanism path is only meaningful if the model actually uses it. For each held-out pair
whose bridge gene is curated, we delete the single mechanism edge and re-score, and compare
against deleting a random edge of the same type. Deleting the mechanism edge collapses the
prediction far more than a random edge (paired Wilcoxon $p\approx3.6\times10^{-34}$; faithful
in $83.3\%$ of cases). The mechanism the system names is the mechanism the prediction leans
on, not a post-hoc story.

We also test mechanism \emph{recovery}: with the direct drug--gene edge hidden, only the
\emph{jointly-supervised} GNN names the curated bridge gene (R@10 $0.25$), where a
degree-trivial baseline and a link-only GNN scored through the same head recover $\sim$0.
Because the link-only head was never trained on the mechanism task, this isolates the
\emph{objective}, not the architecture: joint mechanism supervision creates a recoverable
mechanism signal that a link objective alone does not. We do not claim de novo discovery;
this is recovery of \emph{known} mechanism through indirect structure on a modest sample.

\section{A popularity-deconfounded orthogonal check (an honest negative)}\label{sec:ct}
Do the model's novel predictions show up as real human trials? A naive check against
ClinicalTrials.gov is badly confounded: some drugs are trialed against every cancer and some
cancers against every drug. We fit an additive two-way model
$\text{score}(d,c)=\mu+\alpha_d+\beta_c+\varepsilon_{dc}$, remove the drug effect $\alpha_d$
and cancer effect $\beta_c$, and test whether the residual interaction still predicts a
trial. The within-drug signal looks strong ($0.821$), but once the cancer effect is also
removed the interaction AUROC is $0.475$ ($p=0.85$) --- no signal. The apparent effect was
cancer popularity, not pairwise specificity. We report this as a fully characterized
negative; it stands in deliberate contrast to the prospective result, where specific signal
\emph{does} appear.

\section{Discussion and limitations}
The strongest, most defensible framing of this project is not ``our GNN beats XGBoost''
(it does not), but: \emph{a tuned tabular model is a strong ranker, the graph earns its place
by producing testable mechanism paths, and an LLM/retrieval layer checks whether those paths
are supported by the literature.} Every claim is paired with the test that could have
falsified it.

Limitations are real and stated throughout. The prospective year proxy is a literature
co-mention, not a regulatory date. The faithfulness and recovery tests concern \emph{known}
mechanism recovered through indirect structure, not de novo discovery, on modest samples.
The mechanism signal is only weakly discriminative against shared-target negatives. The
ClinicalTrials.gov check is a negative once popularity is removed. And the candidate
shortlist is hypothesis-generating, not a vetted clinical recommendation. We view these not
as caveats to bury but as the substance of an honest evaluation.

\begin{thebibliography}{9}
\bibitem{primekg} Chandak, Huang, Zitnik. Building a knowledge graph to enable precision
medicine (PrimeKG). \emph{Scientific Data}, 2023.
\bibitem{txgnn} Huang, Chandak, \emph{et al.}, Zitnik. A foundation model for
clinician-centered drug repurposing (TxGNN). \emph{Nature Medicine}, 2024.
\bibitem{kgmlxdtd} KGML-xDTD: a knowledge graph-based machine learning framework for drug
treatment prediction and mechanism description. arXiv:2212.01384.
\bibitem{decagon} Zitnik, Agrawal, Leskovec. Modeling polypharmacy side effects with graph
convolutional networks (Decagon). \emph{Bioinformatics}, 2018.
\end{thebibliography}

\end{document}
